\section*{Chapitre \ref{ch:b2:alpine-orogen} --- Anatomie d'une chaîne de montagnes~: les Alpes}

\begin{solution}{exo:b2:alpine-orogen:1}
(a) $e = 1 - 27.0/36.0 = 0.25$, soit \qty{25}{\percent}~; $S = 36.0 - 27.0 =
\qty{9.0}{km}$. (b) Le niveau doit avoir conservé sa longueur (ni étirement ni amincissement dans sa longueur),
la coupe doit être parallèle à la direction de transport, et les lignes d'ancrage doivent marquer des lieux
où les couches n'ont pas glissé les unes contre les autres.
\end{solution}

\begin{solution}{exo:b2:alpine-orogen:2}
(a) Une \omterm{def:b2:alpine-orogen:klippe}{klippe}~: le gneiss est un vestige isolé d'une \omterm{def:b2:alpine-orogen:nappes}{nappe de charriage} reposant sur son autochtone~;
l'érosion a enlevé le reste de la feuille tout autour. (b) Une \omterm{def:b2:alpine-orogen:klippe}{fenêtre tectonique}~: l'érosion a
entaillé la \omterm{def:b2:alpine-orogen:nappes}{nappe} de gneiss et mis à nu les roches plus jeunes sous son chevauchement de base, qui
plonge donc en s'éloignant de la fenêtre de tous les côtés.
\end{solution}

\begin{solution}{exo:b2:alpine-orogen:3}
(a) Complexe ultramafique (le manteau supérieur)~; complexe gabbroïque (la croûte océanique inférieure)~;
complexe de \omterm{def:b2:magma-emplacement:bodies}{dykes} en filons (la croûte moyenne)~; complexe volcanique en coussins (la croûte supérieure).
(b) Le complexe de \omterm{def:b2:magma-emplacement:bodies}{dykes} en filons. (c) La plus grande part de la croûte des océans larges est consommée par la subduction~;
les \omterm{def:b2:alpine-orogen:ophiolite}{ophiolites} qui subsistent sur les continents se sont formées surtout dans de petits bassins océaniques ou entre des
arcs insulaires, les \omterm{def:b2:alpine-orogen:ophiolite}{ophiolites} téthysiennes là où un continent s'est aminci jusqu'à l'ouverture d'un petit océan.
\end{solution}

\begin{solution}{exo:b2:alpine-orogen:4}
(a) \omterm{def:b2:alpine-orogen:nappes}{Nappes} du Mont-Rose et du Grand-Saint-Bernard~: le côté européen (le Briançonnais et la
marge passive européenne)~; Zermatt--Saas et zone du Combin~: l'océan Piémont--Ligurie~; Dent
Blanche~: la marge de l'Apulie. (b) L'Apulie, qui se retrouve au sommet. (c) La plaque supérieure n'est pas
subduite, et les écailles raclées de la plaque inférieure s'empilent sous elle dans l'ordre dans lequel la
plaque inférieure est arrivée~: l'océan lointain d'abord, la marge proche en dernier et la plus basse. Un plissement ultérieur de l'empilement,
ou des chevauchements hors séquence, peut bouleverser cet ordre.
\end{solution}

\begin{solution}{exo:b2:alpine-orogen:5}
(a) $2\pi G\,\Delta\rho = 2\pi\times6.674\times10^{-11}\times500 =
\qty{2.097e-7}{s^{-2}}$~; $h = 1.70\times10^{-3}/2.097\times10^{-7} = \qty{8.1}{km}$.
(b) $2.097\times10^{-7}\times25\,000 = \qty{5.24e-3}{m.s^{-2}}$, environ
\qty{-524}{mGal}. (c) La racine n'est pas une plaque infinie~: c'est un corps de quelques dizaines de kilomètres
de profondeur et à peine aussi large que la chaîne, si bien qu'il tire moins sur un point de la surface qu'une
plaque de même épaisseur s'étendant à l'infini~; les corps denses comme la zone d'Ivrée ajoutent des
anomalies positives, et le contraste de densité en profondeur peut être plus faible que supposé.
\end{solution}

\begin{solution}{exo:b2:alpine-orogen:6}
(a) $164 - 44 = \qty{120}{Ma}$. (b) $z = P/(\bar\rho g)$~:
$2.8\times10^9/(3000\times9.8) = \qty{95.2}{km}$; $2.8\times10^9/(3300\times9.8) =
\qty{86.6}{km}$. (c) «~Plus de \qty{90}{km}~» se situe dans cet intervalle~: la profondeur dépend de
la densité moyenne choisie pour les terrains sus-jacents autant que de la pression.
\end{solution}

\begin{solution}{exo:b2:alpine-orogen:7}
(a) Un flanc de longueur $L$ plongeant à $\theta$ s'étend horizontalement sur $L\cos\theta$~; tous les flancs
ont le même pendage, donc $l = l_0\cos\theta$ et $e = 1 - l/l_0 = 1 - \cos\theta$.
(b) $\cos\theta = 0.70$~: $\theta = \qty{45.6}{\degree}$. (c) $1 - \cos\qty{40}{\degree} =
\qty{23.4}{\percent}$, comme sur la figure en chevrons~; $1 - \cos\qty{60}{\degree}
= \qty{50}{\percent}$.
\end{solution}

\begin{solution}{exo:b2:alpine-orogen:8}
(a) $l_0 = 118/2.0 = \qty{59.0}{km}$~; $S = 59.0 - 42 = \qty{17.0}{km}$~; $e = 17.0/59.0 =
\qty{28.8}{\percent}$. (b) $A_e = 118 - 42\times2.0 = \qty{34.0}{km^2}$~; $d = 34.0/17.0 =
\qty{2.0}{km}$. (c) Oui~: le \omterm{def:b2:alpine-orogen:decollement}{décollement} se situe à la base de la couverture de \qty{2.0}{km}, sous
son sommet~; la couverture est détachée par sa base.
\end{solution}

\begin{solution}{exo:b2:alpine-orogen:9}
(a) La racine d'Airy vaut $4.5h$~: à \qty{46.5}{\degree} N, $4.5\times2.74 = \qty{12.3}{km}$, soit un
Moho à \qty{42.3}{km}~; à \qty{45.55}{\degree} N, $4.5\times0.65 = \qty{2.9}{km}$, soit un Moho
à \qty{32.9}{km}. (b) Sous le relief élevé, le Moho est à \qty{6.0}{km} moins profond que ne le prévoit
Airy~; au bord de la plaine du Pô, il est plus profond de \qty{28.3}{km}. (c) La racine n'est pas
compensée localement~: elle est déplacée vers le sud et plus grande que ne l'exige le relief.
La plaque se courbe dans son ensemble sous les charges qui la pèsent (le soutien régional traité dans un chapitre du volume de troisième
année), et l'une de ces charges, le corps d'Ivrée, n'est pas du tout du relief.
\end{solution}

\begin{solution}{exo:b2:alpine-orogen:10}
(a) $5.1/11.3 = 0.45$~: moins de la moitié. (b) $5.1 + 6.2 = \qty{11.3}{km}$. Le toit du
Jurassique se situe au-dessus du \omterm{def:b2:alpine-orogen:decollement}{décollement} supérieur~; les \qty{6.2}{km} transportés le long de ce \omterm{def:b2:alpine-orogen:decollement}{décollement}
sont passés sous les plis de la coupe et ont été dépensés dans des structures plus au nord, si bien que le
niveau supérieur n'a jamais été raccourci de cette quantité dans la coupe. (c) Un seul niveau repère n'enregistre
que le raccourcissement des couches auxquelles il appartient~; chaque \omterm{def:b2:alpine-orogen:decollement}{décollement} situé au-dessus ou au-dessous de lui peut
transporter du déplacement dans la coupe ou hors de celle-ci. Équilibrer plusieurs niveaux, et vérifier avec les aires.
\end{solution}

\begin{solution}{exo:b2:alpine-orogen:11}
(a) Pour un décrochement de \qty{100}{km}, $\sqrt{200^2 + 100^2} = \qty{223.6}{km}$~; pour
\qty{50}{km}, $\sqrt{200^2 + 50^2} = \qty{206.2}{km}$~; pour \qty{150}{km},
$\sqrt{200^2 + 150^2} = \qty{250.0}{km}$. (b) $\arctan(100/200) = \qty{26.6}{\degree}$ par rapport à
la direction nord--sud, une ligne de convergence proche de N\qty{150}{\degree}E citée par Rosenberg et
ses collègues. (c) Le décrochement est incertain de $\pm\qty{50}{km}$, ce qui déplace la convergence globale d'environ
\qty{20}{km} dans un sens ou l'autre~; le troisième chiffre du résultat est très inférieur à
cette incertitude.
\end{solution}

\begin{solution}{exo:b2:alpine-orogen:12}
(a) $210/12.9 = \qty{16.3}{Ma}$ à $210/8.5 = \qty{24.7}{Ma}$ (\unit{km} divisé par
\unit{mm.an^{-1}} donne des millions d'années). (b) Plus de $4\times60 = \qty{240}{km}$~; au
moins $240 - 60 = \qty{180}{km}$ du bassin sont aujourd'hui sous les chevauchements alpins ou ont été
érodés. (c) Le bassin se forme devant le prisme de chevauchement et migre avec lui~: sa partie interne est
toujours recouverte par les chevauchements qui le chargent, et ses plus anciens remplissages sont pris
dans les chevauchements, comme la \omterm{def:b2:alpine-orogen:foreland}{molasse} subalpine.
\end{solution}

\begin{solution}{pb:b2:alpine-orogen:1}
\noindent\textbf{Partie I.}

\textbf{1.} $3.0/\cos\qty{35}{\degree} = \qty{3.66}{km}$~; $2.5/\cos\qty{50}{\degree} =
\qty{3.89}{km}$~; $4.0/\cos\qty{40}{\degree} = \qty{5.22}{km}$.

\textbf{2.} $8.00 + 2\times3.66 + 5.00 + 2\times3.89 + 6.0/\cos\qty{10}{\degree} +
2\times5.22 + 10.0/\cos\qty{5}{\degree} = 8.00 + 7.32 + 5.00 + 7.78 + 6.09 + 10.44 + 10.04 =
\qty{54.67}{km}$~; avec le palier répété, $l_0 = \qty{60.7}{km}$.

\textbf{3.} Le niveau situé sous l'écaille de chevauchement est le même niveau, entraîné sous l'écaille par le
chevauchement~: avant le chevauchement, il était aligné avec le reste. L'omettre restaurerait le
niveau trop court.

\textbf{4.} $e = 1 - 48.0/60.7 = 0.209$, soit \qty{20.9}{\percent}~; $S = 60.7 - 48.0 =
\qty{12.7}{km}$.

\textbf{5.} $60.7 - 60.2 = \qty{0.5}{km}$, $0.5/60.7 = \qty{0.8}{\percent}$~: dans la
précision de la mesure~; la coupe est acceptable.

\textbf{6.} $60.7 - 54.0 = \qty{6.7}{km}$, \qty{11}{\percent}~: les deux niveaux d'une même écaille ne peuvent
se restaurer à des longueurs aussi différentes. La coupe ne serait pas admissible telle qu'elle est tracée~: une
structure absente ou mal tracée, une ligne d'ancrage mal placée, un mouvement hors du plan de
la coupe, ou un \omterm{def:b2:alpine-orogen:decollement}{décollement} entre les deux niveaux.

\noindent\textbf{Partie II.}

\textbf{7.} $146/2.4 = \qty{60.8}{km}$.

\textbf{8.} Elle diffère de \qty{60.7}{km} de \qty{0.1}{km}~: la longueur et l'aire concordent.

\textbf{9.} $A_e = 146 - 48.0\times2.4 = 146 - 115.2 = \qty{30.8}{km^2}$.

\textbf{10.} $d = 30.8/12.7 = \qty{2.43}{km}$, égale aux \qty{2.4}{km} de la couverture à
la précision des données~: le \omterm{def:b2:alpine-orogen:decollement}{décollement} est la base de la couverture.

\textbf{11.} $160/2.4 = \qty{66.7}{km}$, soit \qty{6.0}{km} de plus que la longueur de ligne. L'épaississement
des évaporites à la base de la couverture, qui s'écoulent dans le cœur des anticlinaux, ajoute de l'aire
sans ajouter de longueur~; l'équilibrage des aires, à épaisseur constante, échoue alors.

\textbf{12.} La couverture s'est déplacée sur un \omterm{def:b2:alpine-orogen:decollement}{décollement} à sa base, et le socle n'a pas été
impliqué. Sous la chaîne, le toit du socle devrait être plat ou pencher doucement et
régulièrement, sans être brisé par les chevauchements qui le surmontent.

\noindent\textbf{Partie III.}

\textbf{13.} $22 + 18 + 25 = \qty{65}{km}$~; $S = 65 - 15 = \qty{50}{km}$.

\textbf{14.} $e = 50/65 = \qty{76.9}{\percent}$, contre \qty{20.9}{\percent} dans la
chaîne d'avant-pays~: les \omterm{def:b2:alpine-orogen:nappes}{nappes} sont bien plus raccourcies.

\textbf{15.} $30 - 22 = \qty{8}{km}$ ($e = \qty{26.7}{\percent}$). Le socle lui-même est
raccourci~: \omterm{def:b2:alpine-orogen:skin}{tectonique de socle}.

\textbf{16.} Chaîne d'avant-pays $12.7/(11 - 4) = \qty{1.81}{km/Ma}$, soit
\qty{1.81}{mm} par an~; \omterm{def:b2:alpine-orogen:nappes}{nappes} $50/(30 - 20) = \qty{5.0}{mm}$ par an.

\textbf{17.} Les \omterm{def:b2:alpine-orogen:nappes}{nappes} (de 30 à \qty{20}{Ma}) se sont déformées avant la chaîne d'avant-pays (de 11 à
\qty{4}{Ma})~: la déformation a migré vers l'extérieur, en direction de l'avant-pays. C'est l'ordre des
Alpes~: le chevauchement de Glaris à l'Oligocène et au Miocène, les massifs externes au Miocène,
le Jura au Miocène supérieur et au Pliocène.

\textbf{18.} $12.7 + 50 + 8 = \qty{70.7}{km}$.

\noindent\textbf{Partie IV.}

\textbf{19.} $\qty{0.5}{cm.an^{-1}}\times\qty{40}{Ma} = 0.005\times4\times10^7 =
\qty{2.0e5}{m} = \qty{200}{km}$.

\textbf{20.} $200 + 18 = \qty{218}{km}$.

\textbf{21.} Seul ce qui est conservé ou déduit est compté~: des parties de l'océan ont pu avoir été
subduites avant \qty{85}{Ma} ou plus vite que la vitesse moyenne, des écailles ophiolitiques ont pu être
érodées ou emportées hors de la coupe, et une convergence oblique à la coupe apporte
du fond océanique que la coupe ne voit pas.

\textbf{22.} Zone interne $60 - 20 = \qty{40}{km}$ ($e = \qty{66.7}{\percent}$)~; rétro-prisme
$70 - 48 = \qty{22}{km}$ ($e = \qty{31.4}{\percent}$).

\textbf{23.} $12.7 + 50 + 8 + 40 + 22 = \qty{132.7}{km}$~; dans la fourchette des estimations
pour les Alpes centrales, environ 125 à \qty{160}{km} de raccourcissement collisionnel (avant d'ajouter
le décrochement).

\textbf{24.} Largeur restaurée $153 + 132.7 = \qty{285.7}{km}$~; croûte restaurée
$285.7\times30 = \qty{8571}{km^2}$~; croûte actuelle $153\times45 = \qty{6885}{km^2}$. Manquant~:
$8571 - 6885 = \qty{1686}{km^2}$, soit \qty{19.7}{\percent} de la croûte restaurée~; elle a été
érodée (et déposée dans les \omterm{def:b2:alpine-orogen:foreland}{bassins d'avant-pays}) ou subduite avec la croûte inférieure dans le
manteau. Sans ces pertes, la croûte aurait en moyenne $8571/153 = \qty{56}{km}$ d'épaisseur.
Le raccourcissement collisionnel total à travers la coupe est
\textbf{\qty[text-series-to-math=true]{133}{km}}.
\end{solution}
